% Fisheries Population Models Template
% Topics: Stock-recruitment, surplus production, MSY, age-structured models, overfishing
% Style: Fisheries management report with computational analysis

\documentclass[a4paper, 11pt]{article}
\usepackage[utf8]{inputenc}
\usepackage[T1]{fontenc}
\usepackage{amsmath, amssymb}
\usepackage{graphicx}
\usepackage{siunitx}
\usepackage{booktabs}
\usepackage{subcaption}
\usepackage[makestderr]{pythontex}

% Theorem environments
\newtheorem{definition}{Definition}[section]
\newtheorem{theorem}{Theorem}[section]
\newtheorem{example}{Example}[section]
\newtheorem{remark}{Remark}[section]

\title{Fisheries Population Models:\\Stock-Recruitment Dynamics and Sustainable Harvesting}
\author{Marine Biology and Fisheries Management Group}
\date{\today}

\begin{document}
\maketitle

\begin{abstract}
This report presents a comprehensive computational analysis of fisheries population models
used in sustainable fisheries management. We examine surplus production models including
the Schaefer and Fox formulations, calculate maximum sustainable yield (MSY) and optimal
fishing effort, analyze stock-recruitment relationships through Beverton-Holt and Ricker
models, and investigate age-structured population dynamics. The analysis demonstrates
critical thresholds for overfishing, equilibrium biomass levels, and yield-per-recruit
optimization for effective fisheries conservation.
\end{abstract}

\section{Introduction}

Fisheries population models provide the scientific foundation for sustainable fisheries
management. These models describe how fish stocks respond to fishing pressure and
environmental variation, enabling managers to set harvest limits that balance economic
productivity with long-term population viability. Understanding stock dynamics is crucial
for preventing overfishing and ensuring food security for coastal communities worldwide.

\begin{definition}[Sustainable Yield]
The sustainable yield is the harvest rate that can be maintained indefinitely without
depleting the stock. Maximum sustainable yield (MSY) represents the largest average catch
that can be taken from a stock under existing environmental conditions.
\end{definition}

\section{Surplus Production Models}

Surplus production models describe population biomass dynamics without explicit age
structure. These models assume that fish populations grow logistically until limited by
carrying capacity, and that fishing removes a portion of the biomass proportional to
fishing effort and stock size.

\subsection{The Schaefer Model}

\begin{definition}[Schaefer Model]
The Schaefer model assumes logistic population growth with fishing mortality proportional
to effort and biomass:
\begin{equation}
\frac{dB}{dt} = rB\left(1 - \frac{B}{K}\right) - qEB
\end{equation}
where $B$ is biomass (tonnes), $r$ is intrinsic growth rate (year$^{-1}$), $K$ is carrying
capacity (tonnes), $E$ is fishing effort (boat-days), and $q$ is catchability coefficient
(boat-days$^{-1}$).
\end{definition}

\begin{theorem}[Maximum Sustainable Yield]
At equilibrium ($dB/dt = 0$), the Schaefer model yields:
\begin{equation}
MSY = \frac{rK}{4}, \quad B_{MSY} = \frac{K}{2}, \quad E_{MSY} = \frac{r}{2q}
\end{equation}
The maximum sustainable yield occurs at half the carrying capacity with fishing effort
equal to half the intrinsic growth rate divided by catchability.
\end{theorem}

\subsection{The Fox Model}

\begin{definition}[Fox Model]
The Fox model assumes Gompertz growth rather than logistic:
\begin{equation}
\frac{dB}{dt} = rB\ln\left(\frac{K}{B}\right) - qEB
\end{equation}
This formulation better describes populations with compensatory mortality at low densities.
\end{definition}

\section{Computational Analysis: Surplus Production}

We implement the Schaefer and Fox models to analyze a hypothetical demersal fish stock
under varying fishing pressure. The analysis demonstrates how fishing effort affects
equilibrium biomass, sustainable yield, and the risk of stock collapse. Understanding
these relationships is essential for setting total allowable catch (TAC) limits in
fisheries management plans.

\begin{pycode}
import numpy as np
import matplotlib.pyplot as plt
from scipy.integrate import odeint
from scipy.optimize import fsolve, minimize_scalar

np.random.seed(42)

# Schaefer model parameters (typical demersal fish stock)
r = 0.8  # intrinsic growth rate (year^-1)
K = 100000  # carrying capacity (tonnes)
q = 0.0001  # catchability coefficient (boat-days^-1)

# Calculate MSY analytically
MSY = r * K / 4
B_MSY = K / 2
E_MSY = r / (2 * q)

print(f"Schaefer Model Parameters:")
print(f"MSY = {MSY:.0f} tonnes/year")
print(f"Biomass at MSY = {B_MSY:.0f} tonnes")
print(f"Optimal effort = {E_MSY:.0f} boat-days/year")

# Schaefer differential equation
def schaefer(B, t, r, K, q, E):
    dBdt = r * B * (1 - B / K) - q * E * B
    return dBdt

# Fox differential equation
def fox(B, t, r, K, q, E):
    if B <= 0:
        return 0
    dBdt = r * B * np.log(K / B) - q * E * B
    return dBdt

# Time array
t = np.linspace(0, 50, 500)

# Different fishing effort levels
effort_levels = np.array([0, 2000, 4000, 6000, 8000, 10000])

# Initial biomass (virgin stock)
B0 = K

# Simulate trajectories
schaefer_trajectories = {}
fox_trajectories = {}

for E in effort_levels:
    B_schaefer = odeint(schaefer, B0, t, args=(r, K, q, E))
    B_fox = odeint(fox, B0, t, args=(r, K, q, E))
    schaefer_trajectories[E] = B_schaefer.flatten()
    fox_trajectories[E] = B_fox.flatten()

# Equilibrium biomass and yield curves
effort_range = np.linspace(0, 12000, 200)
B_eq_schaefer = K * (1 - q * effort_range / r)
B_eq_schaefer[B_eq_schaefer < 0] = 0
yield_schaefer = q * effort_range * B_eq_schaefer

# Fox equilibrium (numerical solution)
def fox_equilibrium(B, E, r, K, q):
    if B <= 0:
        return 1e10
    return r * np.log(K / B) - q * E

B_eq_fox = []
for E in effort_range:
    if E < r / q:
        B_eq_fox.append(fsolve(fox_equilibrium, K/2, args=(E, r, K, q))[0])
    else:
        B_eq_fox.append(0)
B_eq_fox = np.array(B_eq_fox)
yield_fox = q * effort_range * B_eq_fox

# Create first figure: biomass trajectories and equilibrium curves
fig = plt.figure(figsize=(14, 10))

# Plot 1: Schaefer biomass trajectories
ax1 = fig.add_subplot(2, 3, 1)
colors = plt.cm.viridis(np.linspace(0, 0.9, len(effort_levels)))
for i, E in enumerate(effort_levels):
    ax1.plot(t, schaefer_trajectories[E], linewidth=2, color=colors[i],
             label=f'E = {E:.0f}')
ax1.axhline(y=B_MSY, color='red', linestyle='--', linewidth=1.5, alpha=0.7,
            label=f'$B_{{MSY}}$ = {B_MSY:.0f}')
ax1.set_xlabel('Time (years)', fontsize=11)
ax1.set_ylabel('Biomass (tonnes)', fontsize=11)
ax1.set_title('Schaefer Model: Biomass Trajectories', fontsize=12, fontweight='bold')
ax1.legend(fontsize=8, loc='right')
ax1.grid(True, alpha=0.3)
ax1.set_ylim(0, K * 1.1)

# Plot 2: Equilibrium biomass vs effort
ax2 = fig.add_subplot(2, 3, 2)
ax2.plot(effort_range, B_eq_schaefer, 'b-', linewidth=2.5, label='Schaefer')
ax2.plot(effort_range, B_eq_fox, 'g-', linewidth=2.5, label='Fox')
ax2.axhline(y=B_MSY, color='red', linestyle='--', linewidth=1.5, alpha=0.7)
ax2.axvline(x=E_MSY, color='red', linestyle='--', linewidth=1.5, alpha=0.7)
ax2.fill_between(effort_range, 0, B_eq_schaefer, where=(B_eq_schaefer < 0.2 * K),
                  alpha=0.2, color='red', label='Overfished')
ax2.set_xlabel('Fishing Effort (boat-days/year)', fontsize=11)
ax2.set_ylabel('Equilibrium Biomass (tonnes)', fontsize=11)
ax2.set_title('Equilibrium Stock Size', fontsize=12, fontweight='bold')
ax2.legend(fontsize=9)
ax2.grid(True, alpha=0.3)

# Plot 3: Yield vs effort (production curve)
ax3 = fig.add_subplot(2, 3, 3)
ax3.plot(effort_range, yield_schaefer, 'b-', linewidth=2.5, label='Schaefer')
ax3.plot(effort_range, yield_fox, 'g-', linewidth=2.5, label='Fox')
ax3.axhline(y=MSY, color='red', linestyle='--', linewidth=1.5, alpha=0.7,
            label=f'MSY = {MSY:.0f}')
ax3.axvline(x=E_MSY, color='red', linestyle='--', linewidth=1.5, alpha=0.7)
ax3.scatter([E_MSY], [MSY], s=150, c='red', marker='*', zorder=5, edgecolor='black')
ax3.set_xlabel('Fishing Effort (boat-days/year)', fontsize=11)
ax3.set_ylabel('Sustainable Yield (tonnes/year)', fontsize=11)
ax3.set_title('Production Curve', fontsize=12, fontweight='bold')
ax3.legend(fontsize=9)
ax3.grid(True, alpha=0.3)

# Plot 4: Yield vs biomass
ax4 = fig.add_subplot(2, 3, 4)
ax4.plot(B_eq_schaefer, yield_schaefer, 'b-', linewidth=2.5, label='Schaefer')
ax4.plot(B_eq_fox, yield_fox, 'g-', linewidth=2.5, label='Fox')
ax4.scatter([B_MSY], [MSY], s=150, c='red', marker='*', zorder=5, edgecolor='black',
            label='MSY point')
ax4.axvline(x=0.2*K, color='orange', linestyle=':', linewidth=2, alpha=0.7,
            label='Overfished threshold')
ax4.set_xlabel('Stock Biomass (tonnes)', fontsize=11)
ax4.set_ylabel('Sustainable Yield (tonnes/year)', fontsize=11)
ax4.set_title('Yield vs Stock Size', fontsize=12, fontweight='bold')
ax4.legend(fontsize=9)
ax4.grid(True, alpha=0.3)

# Plot 5: Fishing mortality rate
ax5 = fig.add_subplot(2, 3, 5)
F_range = q * effort_range
F_MSY = q * E_MSY
ax5.plot(F_range, yield_schaefer, 'b-', linewidth=2.5)
ax5.axvline(x=F_MSY, color='red', linestyle='--', linewidth=1.5, alpha=0.7,
            label=f'$F_{{MSY}}$ = {F_MSY:.3f}')
ax5.axvline(x=r, color='orange', linestyle=':', linewidth=2, alpha=0.7,
            label=f'$F = r$ (collapse)')
ax5.fill_between(F_range, 0, yield_schaefer, where=(F_range > r),
                  alpha=0.3, color='red', label='Overfishing')
ax5.set_xlabel('Fishing Mortality Rate (year$^{-1}$)', fontsize=11)
ax5.set_ylabel('Sustainable Yield (tonnes/year)', fontsize=11)
ax5.set_title('Yield vs Fishing Mortality', fontsize=12, fontweight='bold')
ax5.legend(fontsize=9)
ax5.grid(True, alpha=0.3)
ax5.set_xlim(0, 1.2)

# Plot 6: Stock depletion levels
ax6 = fig.add_subplot(2, 3, 6)
depletion = B_eq_schaefer / K * 100
ax6.plot(effort_range, depletion, 'b-', linewidth=2.5)
ax6.axhline(y=50, color='red', linestyle='--', linewidth=1.5, alpha=0.7,
            label='50% (MSY)')
ax6.axhline(y=20, color='orange', linestyle='--', linewidth=1.5, alpha=0.7,
            label='20% (overfished)')
ax6.axhline(y=10, color='darkred', linestyle='--', linewidth=1.5, alpha=0.7,
            label='10% (critical)')
ax6.fill_between(effort_range, 0, depletion, where=(depletion < 20),
                  alpha=0.3, color='red')
ax6.set_xlabel('Fishing Effort (boat-days/year)', fontsize=11)
ax6.set_ylabel('Stock Depletion Level (\\%)', fontsize=11)
ax6.set_title('Depletion vs Effort', fontsize=12, fontweight='bold')
ax6.legend(fontsize=8, loc='upper right')
ax6.grid(True, alpha=0.3)
ax6.set_ylim(0, 105)

plt.tight_layout()
plt.savefig('fisheries_models_production.pdf', dpi=150, bbox_inches='tight')
plt.close()
\end{pycode}

\begin{figure}[htbp]
\centering
\includegraphics[width=\textwidth]{fisheries_models_production.pdf}
\caption{Surplus production model analysis showing biomass trajectories under different
fishing effort levels, equilibrium relationships between effort and stock size, and the
classic dome-shaped production curve. The Schaefer model predicts MSY at 50\% of carrying
capacity with effort = \py{f"{E_MSY:.0f}"} boat-days/year. The red shaded regions indicate
overfished conditions where stock biomass falls below 20\% of carrying capacity, risking
recruitment failure and economic collapse of the fishery.}
\label{fig:production}
\end{figure}

The surplus production analysis reveals that fishing effort above \py{f"{E_MSY:.0f}"}
boat-days/year drives the stock below the MSY biomass level, reducing long-term catch.
At extreme effort levels exceeding \py{f"{r/q:.0f}"} boat-days/year, the fishing
mortality rate surpasses the intrinsic growth rate, causing inevitable stock collapse.
This demonstrates the biological and economic irrationality of overfishing.

\section{Stock-Recruitment Relationships}

Stock-recruitment models describe how spawning stock biomass (SSB) determines the
abundance of recruits entering the fishery. These models capture density-dependent
mortality in early life stages and are essential for understanding population resilience.

\subsection{The Beverton-Holt Model}

\begin{definition}[Beverton-Holt Recruitment]
The Beverton-Holt model describes asymptotic recruitment with density dependence:
\begin{equation}
R = \frac{\alpha S}{1 + \beta S}
\end{equation}
where $R$ is recruitment, $S$ is spawning stock biomass, $\alpha$ is maximum recruits
per spawner at low density, and $\beta$ controls density dependence.
\end{definition}

\subsection{The Ricker Model}

\begin{definition}[Ricker Recruitment]
The Ricker model allows for overcompensation at high spawning stock:
\begin{equation}
R = \alpha S e^{-\beta S}
\end{equation}
This formulation can produce a recruitment maximum at intermediate spawning stock levels
due to cannibalism or disease transmission.
\end{definition}

\section{Computational Analysis: Stock-Recruitment}

We analyze stock-recruitment relationships to identify critical thresholds for population
replacement and assess the risk of recruitment overfishing. The comparison of Beverton-Holt
and Ricker models illustrates how different density-dependent mechanisms affect population
stability and optimal spawning stock levels.

\begin{pycode}
# Stock-recruitment parameters
alpha_BH = 50  # recruits per tonne SSB at low density
beta_BH = 0.01  # density dependence (tonne^-1)
alpha_Ricker = 20  # low-density productivity
beta_Ricker = 0.00005  # density dependence

# Spawning stock biomass range
SSB = np.linspace(0, 50000, 500)

# Beverton-Holt recruitment
R_BH = (alpha_BH * SSB) / (1 + beta_BH * SSB)

# Ricker recruitment
R_Ricker = alpha_BH * SSB * np.exp(-beta_Ricker * SSB)

# Replacement line (R = S with survival rate)
survival_rate = 0.05  # fraction of recruits surviving to spawn
replacement = SSB / survival_rate

# Find equilibrium points
def find_equilibria(SSB_range, R_model, survival):
    equil_points = []
    R_vals = R_model(SSB_range) if callable(R_model) else R_model
    for i in range(len(SSB_range) - 1):
        if (R_vals[i] * survival - SSB_range[i]) * (R_vals[i+1] * survival - SSB_range[i+1]) < 0:
            equil_points.append(SSB_range[i])
    return equil_points

# Spawning potential ratio (SPR) analysis
fishing_mortality = np.linspace(0, 2.0, 100)
SPR = np.exp(-fishing_mortality)  # simplified SPR
reference_SPR = np.exp(-0.4)  # F40% reference point

# Calculate recruits per spawning stock (R/S)
R_per_S_BH = R_BH / (SSB + 1e-6)
R_per_S_Ricker = R_Ricker / (SSB + 1e-6)

# Find Ricker optimum
SSB_opt_Ricker = 1 / beta_Ricker
R_opt_Ricker = alpha_BH * SSB_opt_Ricker * np.exp(-1)

# Create second figure: stock-recruitment
fig2 = plt.figure(figsize=(14, 10))

# Plot 1: Stock-recruitment curves
ax1 = fig2.add_subplot(2, 3, 1)
ax1.plot(SSB, R_BH, 'b-', linewidth=2.5, label='Beverton-Holt')
ax1.plot(SSB, R_Ricker, 'g-', linewidth=2.5, label='Ricker')
ax1.plot(SSB, replacement, 'r--', linewidth=2, alpha=0.7, label='Replacement')
ax1.scatter([SSB_opt_Ricker], [R_opt_Ricker], s=150, c='green', marker='o',
            zorder=5, edgecolor='black', label='Ricker optimum')
ax1.set_xlabel('Spawning Stock Biomass (tonnes)', fontsize=11)
ax1.set_ylabel('Recruitment (millions of individuals)', fontsize=11)
ax1.set_title('Stock-Recruitment Relationships', fontsize=12, fontweight='bold')
ax1.legend(fontsize=9)
ax1.grid(True, alpha=0.3)

# Plot 2: Recruits per spawner
ax2 = fig2.add_subplot(2, 3, 2)
ax2.plot(SSB, R_per_S_BH, 'b-', linewidth=2.5, label='Beverton-Holt')
ax2.plot(SSB, R_per_S_Ricker, 'g-', linewidth=2.5, label='Ricker')
ax2.axhline(y=1/survival_rate, color='red', linestyle='--', linewidth=2, alpha=0.7,
            label='Replacement level')
ax2.set_xlabel('Spawning Stock Biomass (tonnes)', fontsize=11)
ax2.set_ylabel('Recruits per Tonne SSB', fontsize=11)
ax2.set_title('Density-Dependent Recruitment', fontsize=12, fontweight='bold')
ax2.legend(fontsize=9)
ax2.grid(True, alpha=0.3)
ax2.set_ylim(0, 60)

# Plot 3: Phase diagram (SSB dynamics)
ax3 = fig2.add_subplot(2, 3, 3)
SSB_next_BH = R_BH * survival_rate
SSB_next_Ricker = R_Ricker * survival_rate
ax3.plot(SSB, SSB_next_BH, 'b-', linewidth=2.5, label='Beverton-Holt')
ax3.plot(SSB, SSB_next_Ricker, 'g-', linewidth=2.5, label='Ricker')
ax3.plot(SSB, SSB, 'r--', linewidth=2, alpha=0.7, label='1:1 line')
ax3.fill_between(SSB, 0, SSB_next_BH, where=(SSB_next_BH > SSB),
                  alpha=0.2, color='green', label='Increasing')
ax3.fill_between(SSB, 0, SSB_next_BH, where=(SSB_next_BH < SSB),
                  alpha=0.2, color='red', label='Decreasing')
ax3.set_xlabel('Current SSB (tonnes)', fontsize=11)
ax3.set_ylabel('Next Generation SSB (tonnes)', fontsize=11)
ax3.set_title('Population Dynamics Phase Diagram', fontsize=12, fontweight='bold')
ax3.legend(fontsize=8, loc='upper left')
ax3.grid(True, alpha=0.3)

# Plot 4: Spawning potential ratio
ax4 = fig2.add_subplot(2, 3, 4)
ax4.plot(fishing_mortality, SPR * 100, 'b-', linewidth=2.5)
ax4.axhline(y=40, color='red', linestyle='--', linewidth=2, alpha=0.7,
            label='F40% reference')
ax4.axhline(y=20, color='orange', linestyle='--', linewidth=2, alpha=0.7,
            label='F20% limit')
ax4.fill_between(fishing_mortality, 0, SPR * 100, where=(SPR * 100 < 20),
                  alpha=0.3, color='red', label='Overfished')
ax4.set_xlabel('Fishing Mortality Rate (year$^{-1}$)', fontsize=11)
ax4.set_ylabel('Spawning Potential Ratio (\\%)', fontsize=11)
ax4.set_title('SPR-Based Reference Points', fontsize=12, fontweight='bold')
ax4.legend(fontsize=9)
ax4.grid(True, alpha=0.3)
ax4.set_xlim(0, 2)

# Plot 5: Recruitment with environmental stochasticity
ax5 = fig2.add_subplot(2, 3, 5)
np.random.seed(123)
n_years = 50
SSB_sim = np.zeros(n_years)
R_sim = np.zeros(n_years)
SSB_sim[0] = 20000
sigma_R = 0.4  # recruitment variability

for year in range(n_years - 1):
    R_mean = (alpha_BH * SSB_sim[year]) / (1 + beta_BH * SSB_sim[year])
    R_sim[year] = R_mean * np.exp(sigma_R * np.random.randn() - 0.5 * sigma_R**2)
    SSB_sim[year + 1] = R_sim[year] * survival_rate * 0.7  # with fishing

ax5.plot(range(n_years), SSB_sim, 'b-', linewidth=2, label='SSB')
ax5_twin = ax5.twinx()
ax5_twin.bar(range(n_years - 1), R_sim[:-1], alpha=0.5, color='green', label='Recruitment')
ax5.axhline(y=B_MSY/5, color='red', linestyle='--', linewidth=1.5, alpha=0.7,
            label='Limit reference')
ax5.set_xlabel('Year', fontsize=11)
ax5.set_ylabel('SSB (tonnes)', color='blue', fontsize=11)
ax5_twin.set_ylabel('Recruitment (millions)', color='green', fontsize=11)
ax5.set_title('Stochastic Recruitment Simulation', fontsize=12, fontweight='bold')
ax5.legend(loc='upper left', fontsize=9)
ax5_twin.legend(loc='upper right', fontsize=9)
ax5.grid(True, alpha=0.3)

# Plot 6: Allee effect threshold
ax6 = fig2.add_subplot(2, 3, 6)
# Modified Ricker with Allee effect
SSB_allee = np.linspace(0, 50000, 500)
S_critical = 5000  # critical depensation threshold
allee_factor = SSB_allee / (SSB_allee + S_critical)
R_allee = alpha_BH * SSB_allee * np.exp(-beta_Ricker * SSB_allee) * allee_factor

ax6.plot(SSB_allee, R_allee * survival_rate, 'purple', linewidth=2.5,
         label='With Allee effect')
ax6.plot(SSB, R_Ricker * survival_rate, 'g--', linewidth=2, alpha=0.6,
         label='Without Allee effect')
ax6.plot(SSB_allee, SSB_allee, 'r--', linewidth=2, alpha=0.7, label='Replacement')
ax6.axvline(x=S_critical, color='orange', linestyle=':', linewidth=2, alpha=0.7,
            label=f'Critical threshold')
ax6.fill_between(SSB_allee, 0, 50000, where=(SSB_allee < S_critical),
                  alpha=0.2, color='red', label='Collapse risk')
ax6.set_xlabel('Spawning Stock Biomass (tonnes)', fontsize=11)
ax6.set_ylabel('Next Generation SSB (tonnes)', fontsize=11)
ax6.set_title('Depensation and Allee Effects', fontsize=12, fontweight='bold')
ax6.legend(fontsize=8, loc='upper left')
ax6.grid(True, alpha=0.3)

plt.tight_layout()
plt.savefig('fisheries_models_recruitment.pdf', dpi=150, bbox_inches='tight')
plt.close()
\end{pycode}

\begin{figure}[htbp]
\centering
\includegraphics[width=\textwidth]{fisheries_models_recruitment.pdf}
\caption{Stock-recruitment analysis demonstrating density-dependent recruitment dynamics
in fish populations. The Beverton-Holt model shows asymptotic recruitment approaching a
maximum, while the Ricker model exhibits overcompensation with peak recruitment at
\py{f"{SSB_opt_Ricker:.0f}"} tonnes of spawning stock. Phase diagrams reveal stable
equilibria and population trajectories under fishing. The spawning potential ratio (SPR)
framework identifies F40\% as a precautionary reference point maintaining recruitment
capacity. Stochastic simulations show recruitment variability's impact on stock dynamics,
while Allee effect analysis demonstrates critical thresholds below \py{f"{S_critical:.0f}"}
tonnes where recruitment failure becomes likely.}
\label{fig:recruitment}
\end{figure}

Stock-recruitment analysis demonstrates that maintaining spawning stock above critical
thresholds is essential for population persistence. The Beverton-Holt formulation suggests
resilience to overfishing, while the Ricker model with Allee effects reveals collapse
risk when spawning biomass falls below \py{f"{S_critical:.0f}"} tonnes.

\section{Age-Structured Models: Yield-Per-Recruit}

Age-structured models explicitly track cohorts through their life history, accounting
for size-dependent growth, natural mortality, and fishing selectivity. Yield-per-recruit
(YPR) analysis optimizes fishing patterns to maximize lifetime catch from each cohort.

\begin{pycode}
# Age-structured model parameters
ages = np.arange(0, 21)  # age classes 0-20 years
L_inf = 80  # asymptotic length (cm)
K_growth = 0.15  # von Bertalanffy growth coefficient
t0 = -1.0  # theoretical age at zero length
a_length_weight = 0.01  # length-weight coefficient
b_length_weight = 3.0  # length-weight exponent
M = 0.2  # natural mortality (year^-1)

# Fishing mortality rates to test
F_values = np.linspace(0, 1.5, 100)

# Selectivity (logistic with 50% selection at age 3)
age_50 = 3.0
selectivity_slope = 1.5
selectivity = 1 / (1 + np.exp(-selectivity_slope * (ages - age_50)))

# von Bertalanffy growth
lengths = L_inf * (1 - np.exp(-K_growth * (ages - t0)))
weights = a_length_weight * lengths**b_length_weight / 1000  # kg

# Calculate yield-per-recruit for different F
YPR = np.zeros(len(F_values))
SPR_values = np.zeros(len(F_values))

for i, F in enumerate(F_values):
    survivors = np.zeros(len(ages))
    survivors[0] = 1.0  # start with 1 recruit

    cumulative_catch = 0
    cumulative_spawners = 0

    for age_idx in range(1, len(ages)):
        F_age = F * selectivity[age_idx]
        Z = M + F_age  # total mortality
        survivors[age_idx] = survivors[age_idx - 1] * np.exp(-Z)
        catch = survivors[age_idx - 1] * (F_age / Z) * (1 - np.exp(-Z))
        cumulative_catch += catch * weights[age_idx]
        cumulative_spawners += survivors[age_idx] * weights[age_idx]

    YPR[i] = cumulative_catch
    SPR_values[i] = cumulative_spawners

# Find F_max and F_0.1
F_max_idx = np.argmax(YPR)
F_max = F_values[F_max_idx]

# F_0.1: fishing mortality where slope is 10% of origin slope
origin_slope = (YPR[1] - YPR[0]) / (F_values[1] - F_values[0])
target_slope = 0.1 * origin_slope
slopes = np.diff(YPR) / np.diff(F_values)
F_01_idx = np.argmin(np.abs(slopes - target_slope))
F_01 = F_values[F_01_idx]

# F_40: fishing mortality giving 40% SPR
SPR_unfished = SPR_values[0]
SPR_40_target = 0.4 * SPR_unfished
F_40_idx = np.argmin(np.abs(SPR_values - SPR_40_target))
F_40 = F_values[F_40_idx]

# Create third figure: age-structured analysis
fig3 = plt.figure(figsize=(14, 10))

# Plot 1: Growth curve
ax1 = fig3.add_subplot(2, 3, 1)
ax1.plot(ages, lengths, 'b-', linewidth=2.5)
ax1.axhline(y=L_inf, color='red', linestyle='--', linewidth=1.5, alpha=0.7,
            label=f'$L_\\infty$ = {L_inf} cm')
ax1.scatter(ages[::2], lengths[::2], s=50, c='blue', edgecolor='black', zorder=5)
ax1.set_xlabel('Age (years)', fontsize=11)
ax1.set_ylabel('Length (cm)', fontsize=11)
ax1.set_title('von Bertalanffy Growth Curve', fontsize=12, fontweight='bold')
ax1.legend(fontsize=9)
ax1.grid(True, alpha=0.3)

# Plot 2: Weight-at-age
ax2 = fig3.add_subplot(2, 3, 2)
ax2.plot(ages, weights, 'g-', linewidth=2.5)
ax2.scatter(ages[::2], weights[::2], s=50, c='green', edgecolor='black', zorder=5)
ax2.set_xlabel('Age (years)', fontsize=11)
ax2.set_ylabel('Weight (kg)', fontsize=11)
ax2.set_title('Weight-at-Age', fontsize=12, fontweight='bold')
ax2.grid(True, alpha=0.3)

# Plot 3: Selectivity curve
ax3 = fig3.add_subplot(2, 3, 3)
ax3.plot(ages, selectivity, 'purple', linewidth=2.5)
ax3.axhline(y=0.5, color='gray', linestyle=':', linewidth=1.5, alpha=0.7)
ax3.axvline(x=age_50, color='gray', linestyle=':', linewidth=1.5, alpha=0.7,
            label=f'$a_{{50}}$ = {age_50} years')
ax3.set_xlabel('Age (years)', fontsize=11)
ax3.set_ylabel('Selectivity', fontsize=11)
ax3.set_title('Fishing Gear Selectivity', fontsize=12, fontweight='bold')
ax3.legend(fontsize=9)
ax3.grid(True, alpha=0.3)
ax3.set_ylim(0, 1.1)

# Plot 4: Yield-per-recruit curve
ax4 = fig3.add_subplot(2, 3, 4)
ax4.plot(F_values, YPR, 'b-', linewidth=2.5)
ax4.scatter([F_max], [YPR[F_max_idx]], s=200, c='red', marker='*', zorder=5,
            edgecolor='black', label=f'$F_{{max}}$ = {F_max:.3f}')
ax4.scatter([F_01], [YPR[F_01_idx]], s=150, c='orange', marker='o', zorder=5,
            edgecolor='black', label=f'$F_{{0.1}}$ = {F_01:.3f}')
ax4.axvline(x=F_max, color='red', linestyle='--', linewidth=1.5, alpha=0.5)
ax4.axvline(x=F_01, color='orange', linestyle='--', linewidth=1.5, alpha=0.5)
ax4.axvline(x=F_40, color='green', linestyle='--', linewidth=1.5, alpha=0.5,
            label=f'$F_{{40}}$ = {F_40:.3f}')
ax4.set_xlabel('Fishing Mortality (year$^{-1}$)', fontsize=11)
ax4.set_ylabel('Yield-per-Recruit (kg)', fontsize=11)
ax4.set_title('Yield-per-Recruit Analysis', fontsize=12, fontweight='bold')
ax4.legend(fontsize=9, loc='upper right')
ax4.grid(True, alpha=0.3)

# Plot 5: Spawning potential ratio
ax5 = fig3.add_subplot(2, 3, 5)
ax5.plot(F_values, SPR_values / SPR_unfished * 100, 'g-', linewidth=2.5)
ax5.axhline(y=40, color='red', linestyle='--', linewidth=2, alpha=0.7,
            label='40% SPR reference')
ax5.axhline(y=20, color='orange', linestyle='--', linewidth=2, alpha=0.7,
            label='20% SPR limit')
ax5.axvline(x=F_40, color='green', linestyle='--', linewidth=1.5, alpha=0.5)
ax5.fill_between(F_values, 0, SPR_values / SPR_unfished * 100,
                  where=(SPR_values / SPR_unfished * 100 < 20),
                  alpha=0.3, color='red', label='Overfished')
ax5.set_xlabel('Fishing Mortality (year$^{-1}$)', fontsize=11)
ax5.set_ylabel('Spawning Potential Ratio (\\%)', fontsize=11)
ax5.set_title('SPR vs Fishing Mortality', fontsize=12, fontweight='bold')
ax5.legend(fontsize=9)
ax5.grid(True, alpha=0.3)
ax5.set_xlim(0, 1.5)

# Plot 6: Cohort survival
ax6 = fig3.add_subplot(2, 3, 6)
F_test_values = [0, 0.2, 0.5, 1.0]
colors_cohort = plt.cm.plasma(np.linspace(0.2, 0.9, len(F_test_values)))

for i, F in enumerate(F_test_values):
    survivors_cohort = np.zeros(len(ages))
    survivors_cohort[0] = 1000  # cohort size
    for age_idx in range(1, len(ages)):
        F_age = F * selectivity[age_idx]
        Z = M + F_age
        survivors_cohort[age_idx] = survivors_cohort[age_idx - 1] * np.exp(-Z)
    ax6.plot(ages, survivors_cohort, linewidth=2.5, color=colors_cohort[i],
             label=f'F = {F:.2f}')

ax6.set_xlabel('Age (years)', fontsize=11)
ax6.set_ylabel('Cohort Survivors', fontsize=11)
ax6.set_title('Cohort Survival Curves', fontsize=12, fontweight='bold')
ax6.legend(fontsize=9)
ax6.grid(True, alpha=0.3)
ax6.set_yscale('log')

plt.tight_layout()
plt.savefig('fisheries_models_age_structured.pdf', dpi=150, bbox_inches='tight')
plt.close()
\end{pycode}

\begin{figure}[htbp]
\centering
\includegraphics[width=\textwidth]{fisheries_models_age_structured.pdf}
\caption{Age-structured fisheries analysis incorporating von Bertalanffy growth,
weight-at-age relationships, and logistic gear selectivity. Yield-per-recruit (YPR)
analysis identifies optimal fishing mortality at F\textsubscript{max} = \py{f"{F_max:.3f}"}
year\textsuperscript{-1}, though the more precautionary F\textsubscript{0.1} reference
point of \py{f"{F_01:.3f}"} is recommended to account for recruitment uncertainty.
Spawning potential ratio (SPR) analysis shows that F\textsubscript{40} = \py{f"{F_40:.3f}"}
maintains 40\% of unfished spawning biomass. Cohort survival curves demonstrate the
exponential decay of abundance with age under different fishing pressure levels, with
natural mortality M = \py{f"{M:.2f}"} year\textsuperscript{-1} representing baseline
mortality from predation, disease, and senescence.}
\label{fig:age_structured}
\end{figure}

Yield-per-recruit analysis demonstrates that fishing mortality of \py{f"{F_max:.3f}"}
year\textsuperscript{-1} maximizes catch per cohort, but the F\textsubscript{0.1}
reference point of \py{f"{F_01:.3f}"} provides a more sustainable target that maintains
spawning potential and accounts for recruitment variability in fluctuating environments.

\section{Results Summary}

\begin{pycode}
print(r"\begin{table}[htbp]")
print(r"\centering")
print(r"\caption{Key Fisheries Reference Points and Model Parameters}")
print(r"\begin{tabular}{llc}")
print(r"\toprule")
print(r"Model & Parameter & Value \\")
print(r"\midrule")
print(f"Schaefer & MSY & {MSY:.0f} tonnes/year \\\\")
print(f" & Biomass at MSY & {B_MSY:.0f} tonnes \\\\")
print(f" & Optimal effort & {E_MSY:.0f} boat-days/year \\\\")
print(f" & Fishing mortality at MSY & {q * E_MSY:.3f} year$^{{-1}}$ \\\\")
print(r"\midrule")
print(f"Stock-Recruitment & Beverton-Holt $\\alpha$ & {alpha_BH:.1f} recruits/tonne \\\\")
print(f" & Critical SSB threshold & {S_critical:.0f} tonnes \\\\")
print(f" & Ricker optimum SSB & {SSB_opt_Ricker:.0f} tonnes \\\\")
print(r"\midrule")
print(f"Age-Structured & $F_{{max}}$ & {F_max:.3f} year$^{{-1}}$ \\\\")
print(f" & $F_{{0.1}}$ & {F_01:.3f} year$^{{-1}}$ \\\\")
print(f" & $F_{{40\\%}}$ (SPR) & {F_40:.3f} year$^{{-1}}$ \\\\")
print(f" & Age at 50\\% selectivity & {age_50:.1f} years \\\\")
print(f" & Natural mortality & {M:.2f} year$^{{-1}}$ \\\\")
print(f" & Asymptotic length & {L_inf:.0f} cm \\\\")
print(r"\bottomrule")
print(r"\end{tabular}")
print(r"\label{tab:reference_points}")
print(r"\end{table}")
\end{pycode}

\section{Discussion and Management Implications}

\begin{remark}[Precautionary Approach]
Modern fisheries management adopts precautionary reference points that account for
scientific uncertainty. While F\textsubscript{max} maximizes yield-per-recruit, the
F\textsubscript{0.1} and F\textsubscript{40\%} reference points provide buffers against
recruitment failure and environmental variability. The comparison reveals that sustainable
fishing mortality should be approximately \py{f"{F_01:.3f}"} year\textsuperscript{-1},
substantially below levels that would maximize short-term catch.
\end{remark}

\begin{example}[Rebuilding Overfished Stocks]
For a stock depleted to 15\% of carrying capacity (below the overfished threshold of 20\%),
managers must reduce fishing effort from current levels to allow biomass recovery. Setting
effort below \py{f"{E_MSY/2:.0f}"} boat-days/year creates positive population growth,
enabling stock rebuilding to MSY biomass levels within 10-15 years depending on
recruitment strength.
\end{example}

\subsection{Climate Change Considerations}

Climate-driven changes in ocean temperature and productivity affect all model parameters.
Warming waters may increase growth rates (higher $r$ and $K$) for some species while
reducing carrying capacity for others. Stock-recruitment relationships become more
variable as environmental stochasticity increases, necessitating lower target fishing
mortality rates to maintain resilience.

\subsection{Ecosystem-Based Fisheries Management}

Single-species models provide essential guidance but ignore predator-prey interactions,
habitat degradation, and multispecies harvesting effects. Integrated ecosystem models
extending these frameworks account for trophic dynamics, making harvest recommendations
that sustain ecosystem function rather than maximizing yield from individual species.

\section{Conclusions}

This comprehensive analysis of fisheries population models demonstrates:

\begin{enumerate}
\item The Schaefer surplus production model predicts maximum sustainable yield of
\py{f"{MSY:.0f}"} tonnes/year at optimal fishing effort of \py{f"{E_MSY:.0f}"}
boat-days/year, with equilibrium biomass at 50\% of carrying capacity

\item Stock-recruitment analysis reveals critical spawning biomass thresholds below which
recruitment fails, with the Allee effect creating a collapse threshold at
\py{f"{S_critical:.0f}"} tonnes for this stock

\item Yield-per-recruit analysis identifies F\textsubscript{0.1} = \py{f"{F_01:.3f}"}
year\textsuperscript{-1} as the precautionary fishing mortality target, maintaining
spawning potential while approaching maximum catch efficiency

\item Age-structured models incorporating growth, selectivity, and natural mortality
provide more realistic management advice than biomass-only models, particularly for
long-lived species with delayed maturity
\end{enumerate}

Sustainable fisheries management requires maintaining spawning stock above critical
thresholds, limiting fishing mortality to precautionary levels below F\textsubscript{MSY},
and adapting to environmental change through regular stock assessments. The models
presented here form the scientific foundation for Total Allowable Catch (TAC) limits,
effort restrictions, and rebuilding plans that sustain both fish populations and fishing
communities.

\section*{References}

\begin{itemize}
\item Schaefer, M. B. (1954). Some aspects of the dynamics of populations important to
the management of commercial marine fisheries. \textit{Bulletin of the Inter-American
Tropical Tuna Commission}, 1(2), 27-56.

\item Beverton, R. J. H., \& Holt, S. J. (1957). \textit{On the Dynamics of Exploited
Fish Populations}. Chapman and Hall, London.

\item Ricker, W. E. (1954). Stock and recruitment. \textit{Journal of the Fisheries
Research Board of Canada}, 11(5), 559-623.

\item Hilborn, R., \& Walters, C. J. (1992). \textit{Quantitative Fisheries Stock
Assessment: Choice, Dynamics and Uncertainty}. Chapman and Hall, New York.

\item Quinn, T. J., \& Deriso, R. B. (1999). \textit{Quantitative Fish Dynamics}.
Oxford University Press, Oxford.

\item Fox, W. W. (1970). An exponential surplus-yield model for optimizing exploited
fish populations. \textit{Transactions of the American Fisheries Society}, 99(1), 80-88.

\item Clark, C. W. (1990). \textit{Mathematical Bioeconomics: The Optimal Management
of Renewable Resources}, 2nd ed. Wiley-Interscience, New York.

\item Haddon, M. (2011). \textit{Modelling and Quantitative Methods in Fisheries}, 2nd ed.
Chapman \& Hall/CRC Press, Boca Raton.

\item Walters, C., \& Martell, S. J. D. (2004). \textit{Fisheries Ecology and Management}.
Princeton University Press, Princeton.

\item Mace, P. M., \& Sissenwine, M. P. (1993). How much spawning per recruit is enough?
\textit{Canadian Special Publication of Fisheries and Aquatic Sciences}, 120, 101-118.

\item Punt, A. E., \& Smith, A. D. M. (1999). Harvest strategy evaluation for the
eastern stock of gemfish (\textit{Rexea solandri}). \textit{ICES Journal of Marine
Science}, 56(6), 860-875.

\item Goodyear, C. P. (1993). Spawning stock biomass per recruit in fisheries management:
foundation and current use. \textit{Canadian Special Publication of Fisheries and Aquatic
Sciences}, 120, 67-81.

\item Myers, R. A., Barrowman, N. J., Hutchings, J. A., \& Rosenberg, A. A. (1995).
Population dynamics of exploited fish stocks at low population levels. \textit{Science},
269(5227), 1106-1108.

\item Hilborn, R. (2007). Managing fisheries is managing people: what has been learned?
\textit{Fish and Fisheries}, 8(4), 285-296.

\item Pauly, D., Christensen, V., Guénette, S., et al. (2002). Towards sustainability
in world fisheries. \textit{Nature}, 418(6898), 689-695.

\item Worm, B., Hilborn, R., Baum, J. K., et al. (2009). Rebuilding global fisheries.
\textit{Science}, 325(5940), 578-585.

\item Costello, C., Ovando, D., Hilborn, R., et al. (2012). Status and solutions for
the world's unassessed fisheries. \textit{Science}, 338(6106), 517-520.

\item Thorson, J. T., Minto, C., Minte-Vera, C. V., Kleisner, K. M., \& Longo, C. (2013).
A new role for effort dynamics in the theory of harvested populations and data-poor stock
assessment. \textit{Canadian Journal of Fisheries and Aquatic Sciences}, 70(12), 1829-1844.

\item Free, C. M., Thorson, J. T., Pinsky, M. L., et al. (2019). Impacts of historical
warming on marine fisheries production. \textit{Science}, 363(6430), 979-983.

\item Froese, R., Winker, H., Coro, G., et al. (2018). Status and rebuilding of European
fisheries. \textit{Marine Policy}, 93, 159-170.
\end{itemize}

\end{document}
